\clearpage
\section{Étape II \textendash{} Sélectionner un projet}

\subsection{Synthèse du projet définitif}

\begin{center}
\begin{tabularx}{\textwidth}{P{3.6cm}L}
\toprule
\thead{Entreprise} & \nomentreprise{}, jeune entreprise de conception et de vente
d'équipements dédiés au sommeil. \\
\thead{Produit} & Bouchons d'oreilles à réduction active de bruit conçus
exclusivement pour le sommeil, au format ultra-plat permettant de dormir sur le
côté. Produit autonome, sans Bluetooth ni application. \\
\thead{Problème traité} & Le bruit nocturne fragmente le sommeil. Les bouchons
en mousse classiques atténuent mal les basses fréquences, et les écouteurs
grand public ne sont pas conçus pour être portés une nuit entière. \\
\thead{Clients} & Particuliers adultes en environnement bruyant : logement
urbain, voisinage, conjoint qui ronfle, travail en horaires décalés. \\
\thead{Distribution} & Vente en ligne directe depuis notre propre site, sans
intermédiaire ni réseau de revendeurs. \\
\thead{Revenus} & Vente du produit, complétée par les embouts et accessoires de
remplacement. \\
\thead{Différenciation} & Le confort et la tenue de l'embout, axe sur lequel les
tests indépendants relèvent des faiblesses chez les produits concurrents. \\
\bottomrule
\end{tabularx}
\end{center}

\textbf{Positionnement.} Là où les produits existants dérivent d'écouteurs
adaptés au coucher, notre produit part du besoin inverse : ne rien diffuser,
seulement supprimer le bruit.

\textbf{Un besoin établi.} Selon l'enquête de l'Institut National du Sommeil et
de la Vigilance de mars 2026, 36 \% des Français se disent gênés par le bruit
la nuit, avec 1,8 réveil nocturne en moyenne, et 42 \% placent le fait de bien
dormir en tête des piliers de la santé (INSV, 2026).

\textbf{Un marché solvable.} Les produits concurrents se vendent entre 180 et
260 euros. Il ne s'agit donc pas de créer un besoin, mais de mieux répondre à
un besoin déjà exprimé et déjà solvable.

\textbf{Phasage.} L'année 1 est consacrée à la conception, au prototypage et à
la validation. La production en série et la commercialisation débutent en année
2. Ce séquencement conditionne le plan de financement de l'étape IX.

\clearpage

\subsection{Argumentation du choix}

\subsubsection{Un besoin massif et mesurable}

Le bruit n'est pas une gêne marginale. L'étude conduite par l'ADEME et le
Conseil National du Bruit chiffre le coût social du bruit en France à 147,1
milliards d'euros par an, en retenant explicitement les perturbations du
sommeil parmi les effets sanitaires évalués (ADEME, 2021). À l'échelle individuelle, l'enquête
de l'INSV citée plus haut confirme que le bruit constitue aujourd'hui l'un des
principaux perturbateurs du sommeil des Français.

Ce point a pesé lourd dans notre décision : parmi les quatre idées comparées à
l'étape I, celle-ci est la seule dont le problème est documenté par des sources
publiques indépendantes, ce qui sécurise la suite de l'étude.

\subsubsection{Un client directement accessible}

C'est le critère qui a fait basculer le classement de faisabilité. Les trois
autres projets supposaient de vendre à un exploitant de salle de sport, à un
gestionnaire d'immeuble ou à un constructeur automobile. Chacun de ces canaux
implique un cycle de vente long, des appels d'offres ou des contrats-cadres
inaccessibles à une entreprise naissante.

Notre produit se vend au contraire directement au particulier, en ligne. Nous
maîtrisons la relation client de bout en bout, nous encaissons sans
intermédiaire, et nous pouvons lancer l'activité sans réseau commercial.

\subsubsection{Une différenciation identifiée et non inventée}

Le marché est occupé, et nous ne le contestons pas. Notre positionnement ne
repose donc pas sur la nouveauté du concept, mais sur un défaut récurrent des
produits existants.

Les tests indépendants font ressortir un compromis que personne ne résout
entièrement. Le QuietOn 4 est salué pour l'efficacité de sa réduction active et
son autonomie, mais son prix de 259 euros sur le site officiel et l'absence de
toute fonction audio lui sont reprochés (Le Café du Geek, 2025). Sur la génération précédente, la
rédaction de CNN Underscored, tout en jugeant la réduction de bruit efficace,
indique avoir eu des difficultés à obtenir un maintien confortable
(CNN Underscored, 2026). À l'inverse,
le Soundcore Sleep A30, plus complet et moins cher, voit son autonomie tomber à
environ six heures lorsque la réduction active est engagée (Idealo, 2026).

Nous en tirons trois axes de différenciation, qui structureront le plan de
marchéage de l'étape IV :

\begin{itemize}
  \item \textbf{Le maintien et le confort} : proposer plusieurs tailles et
        matériaux d'embouts, et faire de l'ajustement le c\oe{}ur du produit
        plutôt qu'un accessoire livré en supplément.
  \item \textbf{L'autonomie réelle} : garantir une nuit complète avec la
        réduction active engagée, et non en la désactivant.
  \item \textbf{La simplicité} : un produit entièrement autonome, sans
        Bluetooth, sans appairage et sans application. L'utilisateur met les
        bouchons et dort.
\end{itemize}

\subsubsection{Le choix d'un produit sans Bluetooth}

Ce troisième axe résulte d'une décision technique prise en connaissance de
cause, et non d'une simple économie de fonctionnalités. Il convient de
distinguer deux phénomènes physiques souvent confondus. La réduction active de
bruit repose sur des ondes acoustiques : un microphone capte le bruit ambiant,
un processeur calcule l'onde en opposition de phase, un haut-parleur l'émet, et
les deux ondes s'annulent au niveau du tympan. Le Bluetooth relève au contraire
des ondes électromagnétiques radiofréquences. Un dispositif à réduction active
dépourvu de module radio n'émet donc aucun rayonnement électromagnétique.

Ce choix emporte quatre conséquences favorables au projet :

\begin{itemize}
  \item \textbf{Une certification allégée.} En l'absence de module radio, la
        directive européenne relative aux équipements radioélectriques ne
        s'applique pas. Seul le marquage CE demeure requis, ce qui réduit le
        coût et le délai traités à l'étape VIII.
  \item \textbf{Une autonomie supérieure.} La suppression de la liaison radio
        libère de l'énergie et sert directement notre deuxième axe de
        différenciation, la nuit complète avec réduction active engagée.
  \item \textbf{Aucune donnée personnelle collectée.} Sans application ni
        compte utilisateur, le produit ne traite aucune donnée, ce qui écarte
        les obligations liées au RGPD.
  \item \textbf{Une conception simplifiée.} Moins de composants et moins de
        logiciel signifient un développement plus court et un produit plus
        fiable, ce qui compte pour une première série.
\end{itemize}

\begin{hypothese}
Ce choix sera communiqué sur le registre de la simplicité, de l'autonomie et de
l'absence de collecte de données, jamais sur celui de la santé. L'OMS considère
que les connaissances scientifiques actuelles ne confirment pas d'effets
sanitaires liés à une exposition aux champs électromagnétiques de faible
intensité, et l'ANSES, tout en recommandant la précaution, n'identifie pas de
risque avéré (ANSES, 2022). Toute allégation contraire relèverait de la pratique
commerciale trompeuse.
\end{hypothese}

\begin{hypothese}
Ces trois axes constituent des hypothèses de positionnement issues de la
recherche documentaire. Ils devront être confrontés aux réponses de l'étude de
marché conduite à l'étape V avant d'être retenus définitivement.
\end{hypothese}

\subsubsection{Une cohérence avec nos compétences et nos contraintes}

Le projet mobilise des compétences que nous possédons déjà pour la partie
logicielle et la vente en ligne. La partie électronique et acoustique constitue
en revanche notre principale limite, et nous l'assumons : elle sera traitée par
recours à un partenaire de conception, dont le coût est intégré aux besoins
matériels de l'étape VIII.

Enfin, le produit ne dépend d'aucune spécificité locale. Il peut être expédié
partout, ce qui ouvre un marché international sans investissement supplémentaire
de distribution.

\subsubsection{Les risques que nous acceptons}

Par honnêteté méthodologique, nous listons les faiblesses que ce choix nous fait
assumer, plutôt que de les découvrir en cours d'étude :

\begin{itemize}
  \item Un investissement de départ nettement supérieur à celui d'une activité
        de service, traité à l'étape IX.
  \item Une difficulté technique réelle sur la miniaturisation et sur le
        traitement du ronflement, qui est un son non stationnaire.
  \item Une concurrence disposant de moyens de recherche supérieurs aux nôtres.
  \item L'absence de Bluetooth prive le produit du suivi du sommeil et de la
        diffusion audio, proposés par le concurrent le plus accessible. Nous
        acceptons cette contrepartie au profit de l'autonomie et de la
        simplicité.
  \item Des obligations réglementaires à respecter dès la première série :
        marquage CE, sécurité et transport des batteries lithium.
\end{itemize}

\subsubsection{Documents de recherche}

L'ensemble des sources consultées pour aboutir à cette décision, avec la date de
consultation et l'information recherchée, figure en annexe. Les copies d'écran
correspondantes y sont jointes.

\vspace{0.5em}
\textit{Voir Annexe 3 \textendash{} Documents de recherche ayant conduit au choix
du projet.}
